% Standalone IEEE-format review copy of Member 5's contribution only.
% Member 1 assembles the complete team report in Overleaf.
\documentclass[conference]{IEEEtran}
\usepackage{cite}
\usepackage{url}
\usepackage{booktabs}

\begin{document}
\title{Member 5 Contribution Preview: Memory Subsystem and Full-Processor Verification}
\author{\IEEEauthorblockN{Member 5: Balanagu Srikar}}
\maketitle

\begin{center}
\footnotesize Review copy containing only the assigned Member 5 material for the MCA Single-Cycle RV32I Processor project.
\end{center}

\section{Module Design}
% ==============================================================================
% Section: Proposed Solution -> Module Design (Memory Subsystem)
% Owner: Member 5 (Balanagu Srikar) - Memory & Verification
% ==============================================================================

\noindent\textbf{Memory Subsystem (\texttt{imem.v}, \texttt{dmem.v}):} 
The processor implements a modified Harvard memory architecture with isolated instruction and data address spaces, preventing structural memory hazards during single-cycle \texttt{LW} execution. 
Instruction memory (\texttt{imem}) is a 64-word (256-byte) read-only asynchronous array populated at simulation initialization via \texttt{\$readmemh}. Instructions are indexed by discarding the lower two byte-offset bits (\texttt{addr[7:2]}). 
Data memory (\texttt{dmem}) provides 64 words of read/write storage, featuring combinational asynchronous read access for single-cycle load writeback and synchronous clock-edge writing conditioned upon the write enable (\texttt{we}) signal.

The memory design adheres to the RISC-V Unprivileged ISA specification \cite{riscvspec2019} and follows established microarchitecture methodologies outlined by Patterson and Hennessy \cite{patterson2020computer} and Harris and Harris \cite{harris2021digital}. Verilog modeling conforms to IEEE Std 1364-2005 \cite{ieee1364verilog}.

\section{Proposed Solution: Verification Plan}
% ==============================================================================
% Section: Proposed Solution -> Verification Plan
% Owner: Member 5 (Balanagu Srikar) - Memory & Verification
% ==============================================================================

\subsection{Verification Plan}
\label{sec:verification}

Processor verification employed a layered, self-checking simulation methodology executed in Icarus Verilog and visualized using GTKWave. Testing progressed from unit-level module validation to full-datapath architectural verification under realistic machine code execution.

\subsubsection{Hierarchical Test Strategy}
The verification flow comprises two primary phases:
\begin{enumerate}
    \item \textbf{Unit Testbenches:} Independent self-checking harnesses for isolated modules: the Arithmetic Logic Unit (\texttt{tb\_alu.v}), the 32-register file (\texttt{tb\_regfile.v}), and the dual-stage control unit (\texttt{tb\_control.v}). Each unit test bench systematically applied exhaustive or randomized vectors checking boundary values, such as zero-register protection in \texttt{regfile} and arithmetic flags in \texttt{alu}.
    \item \textbf{Integrated System Testbenches:} The top-level testbenches (\texttt{tb\_cpu.v} and \texttt{tb\_cpu\_loop.v}) instantiate the integrated processor datapath alongside instruction and data memory memories, verifying cycle-by-cycle signal propagation, control hazard resolution, and architectural register-memory consistency.
\end{enumerate}

\subsubsection{Program A: Complete ISA Coverage Benchmark}
To guarantee complete coverage of the 10 target instructions within the RV32I base integer specification, \textbf{Program A} executes each instruction format in an interdependent operational sequence:
\begin{itemize}
    \item \textbf{Initialization:} \texttt{ADDI} populates baseline operands ($x_1 = 5, x_2 = 3$).
    \item \textbf{R-Type Arithmetic and Logic:} \texttt{ADD}, \texttt{SUB}, \texttt{AND}, \texttt{OR}, and \texttt{XOR} execute in succession, writing results into destination registers $x_3$ through $x_7$.
    \item \textbf{Data Memory Load/Store Round-Trip:} \texttt{SW} persists the value of $x_3$ (8) to memory address $0\times00$; \texttt{LW} retrieves the stored word into $x_8$, verifying single-cycle read-after-write coherence between memory and register writeback.
    \item \textbf{Conditional Control Flow:} \texttt{BEQ} compares equal registers ($x_3 == x_8 == 8$), taking a $+8$-byte forward branch to bypass an invalidation instruction ($x_9 \leftarrow 1$). Subsequently, \texttt{BNE} evaluates unequal registers ($x_1 \neq x_2$), taking another $+8$-byte forward branch past $x_{10} \leftarrow 1$.
    \item \textbf{Terminal State:} Execution concludes at target register $x_{11} \leftarrow 42$ before entering an idle self-loop (\texttt{BEQ} $x_0, x_0, 0$).
\end{itemize}
The automated harness verifies all architectural registers and memory locations against golden reference values, guaranteeing zero false passes.

\subsubsection{Program B: Iterative Accumulation Benchmark}
\textbf{Program B} computes the sum of the series $\sum_{i=1}^{5} i = 15$ over 19 executed instructions. This test specifically exercises:
\begin{itemize}
    \item Backward branch execution using a negative two's-complement immediate offset ($-8$ bytes via \texttt{BNE}).
    \item Repeated synchronous register writeback and asynchronous dual-port reading across iterative loop states.
    \item Word-aligned data memory storage with non-zero offset addressing (\texttt{SW} $x_1, 4(x_0)$).
\end{itemize}

Waveform dumps in Value Change Dump (\texttt{.vcd}) format were analyzed to inspect control path timing, hazard-free multiplexer selection, and memory read/write synchronization.


\section{Individual Contributions}
\begin{table}[htbp]
\centering
\caption{Member 5 Individual Contribution}
\begin{tabular}{lp{3.1cm}p{3.2cm}}
\toprule
Member & Modules & Report sections \\
\midrule
Member 5 & Instruction Memory (\texttt{imem.v}), Data Memory (\texttt{dmem.v}), Assembly/Hex Test Programs, Full-CPU Testbenches & Verification Plan, Memory Subsystem Design, References, Video Demo Assembly \\
\bottomrule
\end{tabular}
\end{table}

\section{Verification Summary and Artifacts}
All ten required RV32I instructions were verified using self-checking testbenches (\texttt{tb\_cpu.v} and \texttt{tb\_cpu\_loop.v}). Program A verified single-cycle arithmetic, logic, load/store memory round-trip, and forward conditional branch behavior with zero discrepancies. Program B verified 5 iterations of an accumulation loop using backward conditional branching with negative displacement, completing in 19 instructions. Both tests pass automatically in Icarus Verilog with zero errors. VCD traces and GTKWave captures in \texttt{docs/waveforms/cpu.png} provide comprehensive timing and architectural evidence.

\bibliographystyle{IEEEtran}
\bibliography{member5_references}

\end{document}
